\chapter*{Conclusion générale}
\addcontentsline{toc}{chapter}{Conclusion générale}

La réalisation de ce projet de fin d'études a constitué un aboutissement concret de ma formation en \textbf{BTS Développement d'Applications Informatiques (DAI)} à l'établissement Alkendi. Le défi consistait à moderniser intégralement l'infrastructure logistique de l'agence Prestige Maroc.

\textbf{Objectifs atteints et Réussite Technique :} 
J'ai réussi à concevoir, modéliser, et développer un système d'information complet partant de zéro. Le choix de l'architecture \textbf{Microservices} s'est avéré déterminant. En scindant le projet en dépôts distincts (\texttt{enterprise-portal}, \texttt{service-crm}, \texttt{service-location}, \texttt{service-erp}), j'ai pu démontrer la puissance d'une infrastructure découplée, garantissant que la génération lourde de factures n'affecte jamais l'expérience de réservation des clients. Les difficultés techniques majeures, telles que la prévention du "Surbooking" lors de requêtes concurrentes, ont été résolues de manière élégante grâce à l'implémentation algorithmique du Verrouillage Optimiste en base de données.

\textbf{Bilan personnel, Pédagogique et Impact Commercial (ROI) :} 
Sur le plan personnel et académique, ce projet a été une véritable immersion dans le monde professionnel du "Software Engineering". Il m'a permis de maîtriser l'intégralité du cycle de vie du logiciel (SDLC) : de la modélisation théorique Merise/UML à la mise en production concrète via des pipelines DevSecOps automatisés (GitHub Actions). 
Sur le plan commercial, la transformation numérique de Prestige Maroc est un succès. La solution offre un \textbf{Retour sur Investissement (ROI)} mesurable : le temps alloué à la saisie administrative a été réduit de 40\%, le taux d'erreur humaine dans la facturation est tombé à 0\%, et la possibilité de réserver et payer en ligne par carte bancaire 24h/24 et 7j/7 a ouvert l'agence à une nouvelle clientèle internationale, augmentant mécaniquement son chiffre d'affaires.

\textbf{Perspectives d'évolution et Améliorations Futures :}
Bien que l'application livrée soit totalement fonctionnelle et sécurisée en production, une plateforme numérique est un produit vivant voué à évoluer. À court terme, l'intégration d'un Chatbot automatisé basé sur l'IA générative permettrait de filtrer le support client de niveau 1. À moyen terme, l'exploitation des données de réservation stockées dans MongoDB permettrait d'alimenter un modèle de Machine Learning pour faire du \textit{Yield Management} : ajuster dynamiquement les prix de location à la hausse ou à la baisse en fonction des prédictions météorologiques et des saisons touristiques, maximisant ainsi la rentabilité du parc automobile.

\textbf{Veille Technologique, Éthique et Green IT :}
Le métier de développeur nécessitant une remise en question permanente, j'ai maintenu tout au long du développement une veille technologique stricte, m'appuyant sur des plateformes comme GitHub et Dev.to. Par ailleurs, conscient de l'impact de l'informatique sur le changement climatique, j'ai pris des mesures d'éco-conception logicielle (Green IT). L'utilisation de Redis pour mettre les véhicules en cache mémoire, couplée à une API REST légère, divise par 4 l'utilisation du processeur des serveurs, réduisant ainsi la consommation d'électricité et l'empreinte carbone globale du projet. Ce rapport démontre qu'il est possible d'allier haute technologie, rentabilité commerciale, et responsabilité environnementale.
